\documentclass[12pt, letterpaper]{article}

%-------Packages---------
\usepackage{newpxtext,newpxmath} % Palatino-like text & math (modern)
\usepackage{bboldx}
\usepackage{mathtools}
\usepackage{tikz}
\usetikzlibrary{calc,intersections}
\usepackage{tikz-cd}
\usepackage[all,arc]{xy}
\usepackage{graphicx}

\usepackage{enumerate}
\usepackage[dvipsnames]{xcolor}
\usepackage{float}
\usepackage[left=1in,top=1in,right=1in,bottom=1in]{geometry}
\usepackage{fancyhdr}
\usepackage{multicol}
\usepackage{setspace}

\usepackage{dsfont}
\usepackage{mathrsfs}

\usepackage{tcolorbox}
\tcbuselibrary{theorems,skins,breakable}

\usepackage{xparse}
\usepackage{import}
\usepackage[utf8]{inputenc}

\usepackage{hyperref}
\usepackage{cleveref}

\usepackage{xcolor, ifthen, tcolorbox}
\tcbuselibrary{theorems,skins,breakable}
% Theorem System
% The following environments are provided:
%   New Problem:        \begin{problem} ...
%   New Question Header:   \qh (then followed by subproblems)
%   Subproblem:     \begin{subproblem} ...
%   Problem Proof:  \begin{problemproof} ...
%   Proof:          \\begin{pf}{color} ...
%   Remark:         \rmk (sentence), \rmkb (block)

% Define custom colors
\definecolor{thmscol}{HTML}{F4565B} %For Problems
\definecolor{lemscol}{HTML}{FE844C} %For Lemmes
\definecolor{prpscol}{HTML}{00A7EF} %For proposition
\definecolor{corscol}{HTML}{23DAFF} %For corrolaries
\definecolor{rmkscol}{HTML}{6DEEC5} %For remarks
\definecolor{defscol}{HTML}{18C991} %For definitions
\definecolor{exmscol}{HTML}{7DB800} %For examples
\definecolor{clmscol}{HTML}{165c58} %For claims

%Change between dark(black)/light(white) mode
\newcommand{\colormode}{white}
\ifthenelse{\equal{\colormode}{white}}
      {\newcommand{\TextColor}{black}}
      {\newcommand{\TextColor}{white}}
\newcommand{\pbcolormain}{thmscol!80!\colormode}
\newcommand{\pbcolorsecond}{thmscol!38!\colormode}


% ============================
% Problem Counter
% ============================
\newcounter{subproblemcounter}
\newcounter{problemcounter}

\newcommand{\q}{
    \setcounter{subproblemcounter}{0}%
    \stepcounter{problemcounter} % Increment problem counter
    \color{\TextColor}Problem \arabic{problemcounter}: % Display the problem counter
}

\newcommand{\stepsub}{
    \stepcounter{subproblemcounter}%
    \color{\TextColor}(\alph{subproblemcounter})
}
% ============================
% Problem
% ============================
\newtcolorbox{pb}[1]{
  enhanced,
  frame hidden,
  titlerule=0mm,
  toptitle=1mm,
  bottomtitle=1mm,
  fonttitle=\bfseries\large,
  coltitle=black,
  colbacktitle=\pbcolormain,
  colback=\pbcolorsecond,
  title={\q #1},
  coltext=\TextColor
}

\newtcolorbox{spb}[1]{
  enhanced,
  frame hidden,
  titlerule=0mm,
  toptitle=1mm,
  bottomtitle=1mm,
  fonttitle=\bfseries\large,
  coltitle=black,
  colbacktitle=\pbcolormain,
  colback=\pbcolorsecond,
  title={\stepsub #1},
  coltext=\TextColor,
  attach boxed title to top left={xshift=3mm,yshift=-2mm},
  boxed title style={
    colback=\pbcolormain,
    colframe=\pbcolormain,
  }
}

\newtcolorbox{pbh}[1]{
    enhanced,
    frame hidden,
    titlerule=0mm,
    toptitle=1mm,
    bottomtitle=1mm,
    fonttitle=\bfseries\large,
    coltitle=black,
    colbacktitle=\pbcolormain,
    colback=\pbcolorsecond,
    title={\q #1},
    boxrule=0pt,
    interior hidden,
    attach boxed title to top left={xshift=3mm,yshift=-2mm},
    boxed title style={
        colback=\pbcolormain,
        colframe=\pbcolormain,
    }
}

\newenvironment{problem}[1]{%
  \newpage
  \begin{pb}{#1}
    }{
  \end{pb}
}

\newenvironment{subproblem}[1]{%
  \begin{spb}{#1}
    }{
  \end{spb}
}

\newenvironment{problemproof}{%
    \begin{pf}{\pbcolormain}
        }{
    \end{pf}
}

\newcommand{\qh}[1][]{
    \newpage
    \begin{pbh}{#1}\end{pbh} % Display the problem counter
}

% ============================
% Claim
% ============================

\newtcbtheorem[number within=problemcounter]{claim}{Claim}
{
    enhanced,
    frame hidden,
    titlerule=0mm,
    toptitle=1mm,
    bottomtitle=1mm,
    fonttitle=\bfseries\large,
    coltitle=black,
    colbacktitle=clmscol!50!\colormode,
    colback=clmscol!20!\colormode,
}{clm}

\NewDocumentCommand{\clm}{m+m}{
    \begin{claim*}{#1}{}
        #2
    \end{claim*}
}

\newenvironment{clmpf}{
	{\noindent{\it \textbf{Proof.}}
	\setlength{\parindent}{0pt}
	}
	\tcolorbox[blanker,breakable,left=5mm,parbox=false,
    before upper={\parindent15pt},
    after skip=10pt,
	borderline west={1mm}{0pt}{clmscol!50!\colormode}]
}{
    \textcolor{clmscol!50!\colormode}{\hbox{}\nobreak\hfill$\blacksquare$}
    \endtcolorbox
}

\NewDocumentCommand{\clmp}{m+m+m}{
    \begin{claim*}{#1}{}
        #2
    \end{claim*}

    \begin{clmpf}
        #3
    \end{clmpf}
}


% ============================
% Proof
% ============================

\newenvironment{pf}[1]{%
  \def\pfcolor{#1}% Store the color in a macro
  \noindent{\it \textbf{Proof.}}%
  \tcolorbox[blanker,breakable,left=5mm,parbox=false,%
    before upper={\parindent15pt},%
    after skip=10pt,%
    borderline west={1mm}{0pt}{\pfcolor},%
    coltext=\TextColor
  ]%
  \setlength{\parindent}{0pt}%
}{%
  \textcolor{\pfcolor}{\hbox{}\nobreak\hfill$\blacksquare$}%
  \endtcolorbox%
}

\newtcolorbox{solution}{
  blanker,
  breakable,
  left=5mm,
  parbox=false,%
  before upper={\parindent15pt},%
  after skip=10pt,%
  borderline west={1mm}{0pt}{\pbcolormain},%
  coltext=\TextColor
}


% ============================
% Remark
% ============================


\NewDocumentCommand{\rmk}{+m}{
    {\it \color{rmkscol!100!white}#1}
}

\newenvironment{remark}{
    \par
    \vspace{5pt}
    \begin{minipage}{\textwidth}
        {\par\noindent{\textbf{Remark.}}}
        \tcolorbox[blanker,breakable,left=5mm,
        before skip=10pt,after skip=10pt,
        borderline west={1mm}{0pt}{rmkscol!80!\colormode},
        coltext=\TextColor
        ]
}{
        \endtcolorbox
    \end{minipage}
    \vspace{5pt}
}

\NewDocumentCommand{\rmkb}{+m}{
    \begin{remark}
        #1
    \end{remark}
}


% Define a new command for problems
\renewcommand{\headrule}{\color{\TextColor}\hrulefill}
\newcommand{\C}{\mathbb{C}}
\newcommand{\R}{\mathbb{R}}
\newcommand{\Q}{\mathbb{Q}}
\newcommand{\Z}{\mathbb{Z}}
\newcommand{\N}{\mathbb{N}}
\renewcommand{\P}{\mathbb{P}}
\newcommand{\F}{\mathbb{F}}
\newcommand{\contr}{\Rightarrow \Leftarrow}
\newcommand{\s}{\(\,\)}
\renewcommand{\t}[1]{\text{#1}}

\newcommand{\skcgen}{{\sf Gen}} 
\newcommand{\skcenc}{{\sf Enc}}
\newcommand{\skcdec}{{\sf Dec}}
\newcommand{\mac}{{\sf MAC}}
\newcommand{\ver}{{\sf Vrfy}}

% Meta data
\setstretch{1.2}

\begin{document}
\color{\TextColor}
\pagecolor{\colormode}
\setlength{\parindent}{0pt}
\pagestyle{fancy}
\fancyhf{}
\fancyhead[LOE]{\color{\TextColor}\slshape{\textbf{Vincent Ferrara}}}
\fancyhead[ROE]{\color{\TextColor}HW2 --- \thepage}
\fancyhead[C]{\color{\TextColor}EECS 475}

\begin{problem}{}
    Suppose that $G : \{0, 1\}^n \to \{0, 1\}^{2n}$ is a secure pseudorandom generator
that expands an $n$-bit random seed into a $2n$-bit pseudorandom string. For each of the
following questions: first, give a YES/NO answer; if the answer is YES, please prove it; if
the answer is NO, please provide a counter-example
\end{problem}

\begin{subproblem}{}
    Let $\wedge$ denote the bit-wise and operation. Is $G((0^{n-1}|| 1) \wedge s)$ a secure pseudorandom
generator that expands a random seed of length $n$ into a pseudorandom string of
length $2n$?
\end{subproblem}
\begin{problemproof}
    No, this is not secure since for $G(0^{n-1}||1 \wedge s)$ for a random $n$-bit string $s$. $G(0^{n-1}||1 \wedge s)$ is either
    $G(0^n)$ or $G(0^{n-1}||1)$.

    Now define a distinguisher $D$ s.t.
    \[D(y)=\begin{cases*}
        1 \quad \t{if } y \in \{G(0^n), G(0^{n-1}||1)\},\\
        0 \quad \t{otherwise.}
    \end{cases*}\]

    If $y \leftarrow G((0^{n-1}|| 1) \wedge U_n)$, then $\P[D(y)=1]=1$

    If $y \leftarrow U_{2n}$, then $\P[D(y)=1] = 2/2^{2n}$

    Thus,

    \[\P[D(G(0^{n-1}|| 1) \wedge U_n)=1] - \P[D(U_{2n})=1] \geq 1-2^{1-2n}\]

    this is not negligible. Therefore, this is not a secure PRG.
\end{problemproof}

\begin{subproblem}{}
    Is $G(\t{first half of } G(s))\,||\,G(\t{second half of } G(s))$ a secure pseudorandom generator that
maps $n$ bits to $4n$ bits?
\end{subproblem}
\begin{problemproof}
    Consider these distributions:
    \begin{itemize}
        \item $X_0 = \{G(\t{first half of } G(s))\,||\,G(\t{second half of } G(s)) \mid s \overset{\$}{\leftarrow} \{0,1\}^n\}$
        \item $X_1 = \{G(\t{first half of } s)\,||\,G(\t{second half of } s) \mid s \overset{\$}{\leftarrow} \{0,1\}^{2n}\}$ 
                   
                   $= \{G(x) || G(y) \mid x,y \overset{\$}{\leftarrow} \{0,1\}^{n}\}$
        \item $X_2 = \{G(s) || x \mid s \overset{\$}{\leftarrow} \{0,1\}^{n} \t{ and } x \overset{\$}{\leftarrow} \{0,1\}^{2n}\}$
        \item $X_3 = \{x || y \mid x,y \overset{\$}{\leftarrow} \{0,1\}^{2n}\} =  \{s \mid s  \overset{\$}{\leftarrow} \{0,1\}^{4n}\}$
    \end{itemize}

    Since we know that $G$ is a secure PRG we know that 
    
    $\{G(s) \mid s \overset{\$}{\leftarrow} \{0,1\}^{n}\} \approx \{s \mid s \overset{\$}{\leftarrow} \{0,1\}^{2n}\}$

    This implies that $X_0 \approx X_1$ since a $2n$-bit string from $G$ or uniformly random looks the same, so taking the first and second half will also look the same.

    Again since we know that 

    $\{G(s) \mid s \overset{\$}{\leftarrow} \{0,1\}^{n}\} \approx \{s \mid s \overset{\$}{\leftarrow} \{0,1\}^{2n}\}$

    This implies that $X_1 \approx X_2$, then again use the same thing, and we get $X_2 \approx X_3$.

    Therefore, by hybrid argument we have shown that $X_0 \approx X_3$, and thus, this is a secure PRG.
\end{problemproof}

\begin{problem}{}
    Determine whether the following $F : \{0, 1\}^{2n} \times \{0, 1\}^n \to \{0, 1\}^{2n}$ is a pseudo-
random function. If yes, prove it; if not, describe and analyze a feasible distinguisher that
has non-negligible advantage in distinguishing $F$ from random function family.
\[F_k (x) := G(x) \oplus k \t{ where } G : \{0, 1\}^n \to \{0, 1\}^{2n} \t{ is a PRG.}\]
\end{problem}
\begin{problemproof}
    No, this is not a secure PRF.

    First choose distinct $x_1,x_2 \in \{0,1\}^n$. Query the oracle twice to get $y_i \leftarrow \mathcal{O}(x_i)$ then define $D$ s.t.
    \[D(y_1,y_2) = \begin{cases}
        1 \quad \t{if } y_1 \oplus y_2 = G(x_1) \oplus G(x_2),\\
        0 \quad \t{otherwise.}
    \end{cases}\]
    Thus, $\P[D(y_1,y_2)=1 \mid y_i = G(x_i)]=1$, but for a truly random function $F$, since $y_1,y_2$ 
    are independent and uniform across $\{0,1\}^{2n}$ this implies that $y_1 \oplus y_2$ is also uniform across $\{0,1\}^{2n}$.
    Therefore, $\P[D(y_1,y_2)=1 \mid y_i = F(x_i)]=2^{-2n}$.

    Therefore, the advantage is $1-2^{-2n}$ which is non negligible.
\end{problemproof}

\begin{problem}{}
    Give a distinguisher that has non-negligible advantage in distinguishing the
following $F' : \{0, 1\}^{2n} \times \{0, 1\}^{2n} \to \{0, 1\}^{2n}$ from a random function family.

\phantom{text}

$F'_{k_1||k_2} (x_1||x_2) := F_{k_1} (x_1)|| F_{k_2} (x_2)$ where $F : \{0, 1\}^n \times \{0, 1\}^n \to \{0, 1\}^n$ is a PRF and
$|x_1| = |x_2|$. Here $k_1$ and $k_2$ are drawn uniformly from the keyspace of $F$.
\end{problem}
\begin{problemproof}
    First choose distinct $x_1,x_2 \in \{0,1\}^n$. Query the oracle three times
    \begin{itemize}
        \item $\mathcal{O}(x_1 || x_1) = y_1 || y_1'$
        \item $\mathcal{O}(x_1 || x_2) = y_2 || y_2'$
        \item $\mathcal{O}(x_2 || x_1) = y_3 || y_3'$
    \end{itemize}
    then define $D$ s.t.
    \[D(y_1||y_1', y_2|| y_2', y_3||y_3') \begin{cases*}
        1 \quad \t{if } y_1=y_2 \t{ and } y_1' = y_3',\\
        0 \quad \t{otherwise}.
    \end{cases*}\]
    Thus, $\P[D(y_1||y_1', y_2|| y_2', y_3||y_3')=1 \mid \mathcal{O}=F']=1$, but for a truly random function $H$, since
    $x_1 || x_1, x_1 || x_2, x_2 || x_1$ are all distinct this implies that the three outputs from $H$ are all independent and uniform over $\{0,1\}^{2n}$.
    Thus, $\P[D(y_1||y_1', y_2|| y_2', y_3||y_3') \mid \mathcal{O}=H]=2^{-n} \cdot 2^{-n}=2^{-2n}$.

    Therefore, the advantage is $1-2^{-2n}$ which is non negligible.
\end{problemproof}

\begin{problem}{}
    Consider two encryption schemes 
$\Pi_1 = (\mathsf{Gen}_1, \mathsf{Enc}_1, \mathsf{Dec}_1)$ and 
$\Pi_2 = (\mathsf{Gen}_2, \mathsf{Enc}_2, \mathsf{Dec}_2)$ 
over the key space $\mathcal{K}$, message space $\mathcal{M}$ and ciphertext $\mathcal{C}$. 
Suppose that $\mathcal{K} = \mathcal{M} = \mathcal{C}$.
Now suppose that at least one of $\Pi_1$ and $\Pi_2$ is 
CPA secure but you do not know which one. 
Construct a CPA secure encryption scheme using $\Pi_1$ and $\Pi_2$.
\\\\
Please describe your construction and prove it correct and CPA secure, assuming that both encryption scheme are correct, and at least
one of the underlying encryption schemes is CPA secure.
Besides the two encryption schemes $\Pi_1$ and 
$\Pi_2$, 
please do not use other cryptographic 
primitives in your construction.
\end{problem}
\begin{problemproof}
    Define $\Pi' = (\mathsf{Gen}', \mathsf{Enc}', \mathsf{Dec}')$ s.t. $\skcgen'=\skcgen_1||\skcgen_2$,
    $\skcenc'(k_1||k_2, m)=\skcenc_1(k_1 \skcenc_2(k_2, m))$, and $\skcdec'(k_1||k_2, c)=\skcdec_2(k_2, \skcdec_1(k_1, c))$, where $k_1||k_2$ is the key given from $\skcgen'$.

    Correctness:

    Given $\skcgen' = k_1||k_2$. Let $m \in \mathcal{M}$.

    $\skcdec'(k_1||k_2, \skcenc'(k_1||k_2, m))=$ 
    
    $\skcdec_2(k_2, \skcdec_1(k_1, \skcenc_1(k_1, \skcenc_2(k_2, m))))=$
    $\skcdec_2(k_2, \skcenc_2(k_2, m))=m$.

    CPA secure:

    Case 1: $\Pi_1$ is CPA secure.

    Assume $D$ has non negligible advantage against $\Pi'$. Define $C$ s.t.
    \[\begin{tikzcd}
	\begin{array}{c} \square \\\Pi_1 \\ k_1 \leftarrow \skcgen_1 \end{array} && \begin{array}{c} C\\ k_2 \overset{}{\leftarrow}\skcgen_2 \end{array} && D \\
	{c \leftarrow \skcenc_1(k_1, \skcenc_2(k_2,m_i)) } && {} && {} \\
	\begin{array}{c} \end{array} && \begin{array}{c} \end{array} && \phantom{D}
	\arrow["{\skcenc_1(k_1, c)}"', shift right=3, from=1-1, to=1-3]
	\arrow["{\skcenc_2(k_2,m)=c}"', shift right=3, from=1-3, to=1-1]
	\arrow["{\skcenc_2(k_2, c)}"', shift right=3, from=1-3, to=1-5]
	\arrow["m"', shift right=3, from=1-5, to=1-3]
	\arrow[shift right=5, from=1-5, to=1-5, loop, in=325, out=35, distance=10mm]
	\arrow["c"', shift right=3, from=2-1, to=2-3]
	\arrow["{\skcenc_2(k_2,m_0), \skcenc_2(k_2,m_1)}"', shift right=3, from=2-3, to=2-1]
	\arrow["c"', shift right=3, from=2-3, to=2-5]
	\arrow["{m_0,m_1}"',shift right=3, from=2-5, to=2-3]
	\arrow["{\skcenc_1(k_1, c)}"', shift right=3, from=3-1, to=3-3]
	\arrow["{\skcenc_2(k_2,m)=c}"', shift right=3, from=3-3, to=3-1]
	\arrow["{\skcenc_1(k_1, c)}"', shift right=3, from=3-3, to=3-5]
	\arrow["m"', shift right=3, from=3-5, to=3-3]
	\arrow[shift right=5, from=3-5, to=3-5, loop, in=325, out=35, distance=10mm]
    \end{tikzcd}\]
    Then output what $D$ outputs. This implies that $C$ distinguishes $\Pi_2$. This implies if $\Pi_2$ is secure then $\Pi'$ is secure.

    Case 2: $\Pi_2$ is CPA secure.

    Assume $D$ has non negligible advantage against $\Pi'$. Define $C$ s.t.
    \[\begin{tikzcd}
	\begin{array}{c} \square \\\Pi_2 \\ k_2 \leftarrow \skcgen_2 \end{array} && \begin{array}{c} C\\ k_1 \leftarrow \skcgen_1 \end{array} && D \\
	{c \leftarrow \skcenc_2(k_2, m_i)} && {} && {} \\
	{} && {} && {\phantom{D}}
	\arrow["{\skcenc_2(k_2, m)=c}"', shift right=3, from=1-1, to=1-3]
	\arrow["m"', shift right=3, from=1-3, to=1-1]
	\arrow["{{\skcenc_1(k_1, c)}}"', shift right=3, from=1-3, to=1-5]
	\arrow["m"', shift right=3, from=1-5, to=1-3]
	\arrow[shift right=5, from=1-5, to=1-5, loop, in=325, out=35, distance=10mm]
	\arrow["c"', shift right=3, from=2-1, to=2-3]
	\arrow["{m_0,m_1}"', shift right=3, from=2-3, to=2-1]
	\arrow["{\skcenc_1(k_1,c)}"', shift right=3, from=2-3, to=2-5]
	\arrow["{{m_0,m_1}}"', shift right=3, from=2-5, to=2-3]
	\arrow["{\skcenc_2(k_2, m)=c}"', shift right=3, from=3-1, to=3-3]
	\arrow["m"', shift right=3, from=3-3, to=3-1]
	\arrow["{\skcenc_1(k_1,c)}"', shift right=3, from=3-3, to=3-5]
	\arrow["m"', shift right=3, from=3-5, to=3-3]
	\arrow[shift right=5, from=3-5, to=3-5, loop, in=325, out=35, distance=10mm]
\end{tikzcd}\]
    Then output what $D$ outputs. This implies that $C$ distinguishes $\Pi_1$. This implies if $\Pi_1$ is secure then $\Pi'$ is secure.

    Therefore, in both cases we know that $\Pi'$ is secure given $\Pi_1$ or $\Pi_2$ is secure.
\end{problemproof}

\begin{problem}{}
    Let $F\colon \{0,1\}^n \times \{0,1\}^n \rightarrow \{0,1\}^n$ be a PRF. 
  Consider the following\\ $\t{MAC}=(\skcgen,{\sf Mac},{\sf Vrfy})$ given by
  \begin{itemize}
    \item $\skcgen(1^n)$: outputs a uniformly random key $k\gets \{0,1\}^n$.
    \item ${\sf Mac}(k,m)$: given a message $m=m_1\|m_2\in\{0,1\}^{2n}$, 
      where $|m_1|=|m_2|=n$, output the tag $t:=F_k(m_1)\| F_k(F_k(m_2))$.
    \item ${\sf Vrfy}(k,m,t)$: check whether ${\sf Mac}(k,m)=t$.
  \end{itemize}
  Is this MAC secure (existential unforgeabilty under chosen message attack)? Prove your answer either by giving a formal reduction
  to $F$ being a PRF or by constructing a forger with non-negligible advantage
  in the security game against MAC. 
\end{problem}
\begin{problemproof}
    This is not secure. First pick $m \in \{0,1\}^{2n}$ s.t. $m=m_1||m_2$ where $m_1,m_2$ are distinct $n$-bit strings.
    Query the oracle two times
    \begin{itemize}
        \item $\mac(m_1||m_1)=F_k(m_1)||F_k(F_k(m_1))=t_1$
        \item $\mac(m_2||m_2)=F_k(m_2)||F_k(F_k(m_2))=t_2$
    \end{itemize}
    then return $(m, F_k(m_1)||F_k(F_k(m_2)))$ we can make this tag since we know the values from the queries before, and since $m$ has never queried before and the tag sent is valid this implies this is not 
    MAC secure.
\end{problemproof}

\begin{problem}{}
    For each of the following hash functions $H$, determine whether it
  is \emph{necessarily} collision resistant. If it is, describe and
  analyze an appropriate reduction. If it's not, describe and analyze
  a feasible attacker against the collision resistance of $H$.
\end{problem}

\begin{subproblem}{}
    $H^{s}(x) = H_1^{s}(x) \oplus H_2^{s}(x)$, where $\{H^s_1\}, \{H^s_2\}$ are both
    collision-resistant hash functions families.
\end{subproblem}
\begin{problemproof}
    False, let $\{H_1^s\}$ be a collision-resistant hash function family and let $\{H_2^s\}=\{H_1^s\}$. These are both collision
    resistant

    Then consider $H^s(x)=H_1^s(x) \oplus H_2^s(x)=H_1^s(x) \oplus H_1^s(x)=0^n$.

    Thus, define an attacker $D$ s.t. $D$ returns $x,y \in \{0,1\}^n$ s.t. $x \neq y$. Then $H^s(x)=0^n=H^s(y)$. Thus, with probability 1 $D$ can find a collision.

    Therefore, $H$ is not necessarily collision resistant.
\end{problemproof}

\begin{subproblem}{}
    $H^s(x) = \widetilde{H}^s(\widetilde{H}^s(x))$, where $\{\widetilde{H}^s\}$ is a collision-resistant hash
    function family (that can support arbitrary length inputs).
\end{subproblem}
\begin{problemproof}
    Assume there exists a adversary $D$ that can produce a valid collision for $H^s$ with non negligible probability.
\[\begin{tikzcd}
	\square && C && D
	\arrow["s", shift left=3, from=1-1, to=1-3]
	\arrow["{}", shift left=3, from=1-3, to=1-1]
	\arrow["s", shift left=3, from=1-3, to=1-5]
	\arrow["{(m,m*)}", shift left=3, from=1-5, to=1-3]
\end{tikzcd}\]
Case 1: $\widetilde{H}^s(m)=\widetilde{H}^s(m*)$ return $m,m*$

Case 2: $\widetilde{H}^s(m)\neq\widetilde{H}^s(m*)$ return $\widetilde{H}^s(m),\widetilde{H}^s(m*)$

Therefore, in both cases we provide a valid collision for $\widetilde{H}^s$ with non negligible probability.

Therefore, this implies if $\{\widetilde{H}^s\}$ is a collision-resistant hash function family then $\{H^s\}$ is also.
\end{problemproof}
\end{document}